%%REPLACE title="The unproved reflection must remain a hypothesis"
\begin{dcomment}{development}{p37:dont on sait qu'il admet un adjoint}{The reflection does not exist}
Grothendieck wants to convert homotopy invariance of a represented functor into a condition on its representing object. For a reflective localization this is an elementary and powerful argument, and Propositions~1 and~2 below are correct formal statements about any reflective localization. The edition questions the right adjoint to $\varphi:\mathrm{Cat}\,\mathrm{hot}\to\mathrm{Hot}$ on which their application rests. It does not exist.

\textbf{Theorem.}\ours
\textit{The functor $\varphi:\mathrm{Cat}\,\mathrm{hot}\to\mathrm{Hot}$ has no right adjoint. More precisely, there is no small category $K$ with bijections $\pi_0\operatorname{Fun}(X,K)\cong[BX,K(\mathbf Z/2,2)]$ natural in the small category $X$.}

\textit{Proof.}
A right adjoint $\psi$ would give such a $K$, namely $\psi$ of the homotopy type $Y=K(\mathbf Z/2,2)$, since morphisms $X\to K$ in $\mathrm{Cat}\,\mathrm{hot}$ are the connected components of $\operatorname{Fun}(X,K)$. Suppose $K$ is given. Taking $X=e$ shows that $\pi_0K=[\mathrm{pt},Y]$ is one point.

Let $\Omega$ be an infinite set of cardinality greater than that of the set of arrows of $K$, and $G$ the group of even permutations of $\Omega$ with finite support. It is the union of the alternating groups of the finite subsets of $\Omega$, so it is simple, and its cardinality is that of $\Omega$. A functor from the one-object category $\mathsf BG$ to $K$ is an object $k$ with a homomorphism $G\to\operatorname{Aut}(k)$. Its kernel is a normal subgroup; it is not trivial because $\operatorname{Aut}(k)$ has fewer elements than $G$; so the homomorphism is trivial. A natural transformation between two such constant functors is a single arrow of $K$. Hence $\operatorname{Fun}(\mathsf BG,K)\cong K$ and $\pi_0\operatorname{Fun}(\mathsf BG,K)$ is one point.

On the other side $[BG,K(\mathbf Z/2,2)]=H^2(G;\mathbf Z/2)$ classifies the central extensions of $G$ by $\mathbf Z/2$, and one of them does not split. Let $V$ be the real vector space with orthonormal basis $(e_\omega)_{\omega\in\Omega}$ and $\mathrm{Cl}(V)$ its Clifford algebra, with $v^2=-|v|^2$. For $\omega\neq\omega'$ the unit vector $\tau=(e_\omega-e_{\omega'})/\sqrt2$ is invertible in $\mathrm{Cl}(V)$ and $v\mapsto-\tau v\tau^{-1}$ is the reflection of $V$ that exchanges $e_\omega$ and $e_{\omega'}$. Let $\widehat G$ be the group of products of an even number of such vectors. Sending such a product to the corresponding product of transpositions is a surjective homomorphism $\widehat G\to G$; its kernel is contained in $\{\pm1\}$, the kernel of the map from the Pin group to the orthogonal group, and $-1=(\tau\tau')^2\in\widehat G$ for two orthogonal such vectors $\tau,\tau'$. So $\widehat G$ is a central extension of $G$ by $\mathbf Z/2$. The element $g=(\omega_1\,\omega_2)(\omega_3\,\omega_4)$ has order two, and its two lifts $\pm\tau\tau'$ have square $-1$; a splitting would lift $g$ to an element of order two. So $H^2(G;\mathbf Z/2)\neq0$, and the two sets cannot be in bijection.\hfill$\square$

The argument shows more than the statement: a category $K$ represents on $\mathrm{Cat}\,\mathrm{hot}$ only what homomorphisms into its automorphism groups and its arrows can see, and a group that is large and simple is invisible to it. The functors $X\mapsto[BX,Y]$ that are representable on all of $\mathrm{Cat}\,\mathrm{hot}$ are those with $Y$ a $1$-type, represented by a groupoid (the proof for a general $Y$, which replaces $G$ by large simple groups with the homology of spheres, is not reproduced here); the fundamental groupoid $X\mapsto\Pi_1X$ is a reflection of $\mathrm{Cat}\,\mathrm{hot}$ onto groupoids modulo natural isomorphism, and it inverts the functors that are bijective on $\pi_0$ and on all $\pi_1$, a larger class than the weak equivalences. Propositions~1 and~2 apply to that reflection.

Thomason's model structure gives a fibrant replacement $K\to RK$ and an equivalence between its homotopy category and $\mathrm{Hot}$ \citep[Theorem~4.9 and Corollary~5.5]{Thomason1980}. It does not identify $\pi_0\operatorname{Fun}(X,RK)$ with $[NX,NK]$ for every $X$: mapping formulas in a model category need a cofibrant source, and the cofibrant objects of that structure are posets. Below, references to an envelope or to a reflective localization of $\mathrm{Cat}\,\mathrm{hot}$ onto homotopy types are to be read as hypotheses that cannot be met; references to Thomason fibrant replacements mean the model-category construction and carry only its conclusions.
\end{dcomment}

%%BLOCK after="A complete instance of the length obstruction"
\begin{dcomment}{development}{p49:Je présume que non - qu'il peut arriver}{The same ind-morphism, two different maps of the unions}
With Thomason fibrancy imposed on every stage and on the unions, the questions of ms.~47--49 still have negative answers. The example is a strict form of the symmetric-group shift used at ms.~52.

\textbf{Proposition.}\ours
\textit{There are a sequence $T_0\subset T_1\subset\cdots$ of full inclusions of finite categories, each $T_n$ Thomason fibrant and isomorphic in $\mathrm{Cat}\,\mathrm{hot}$ to a group, with Thomason-fibrant union $T$, and a functor $s:T\to T$ with $s(T_n)\subseteq T_{n+1}$, such that:}
\begin{enumerate}[label=(\alph*)]
\item \textit{for every $n$, the functor $s|_{T_n}$ and the inclusion $T_n\to T_{n+1}$ are joined by a zigzag of three natural transformations, so $s$ and $\mathrm{id}_T$ induce the same endomorphism of the ind-object $(T_n)$ of $\mathrm{Cat}\,\mathrm{hot}$;}
\item \textit{$Bs$ is not a weak equivalence; in particular $s$ and $\mathrm{id}_T$ are not homotopic;}
\item \textit{there is a weak equivalence $r:T\to\mathsf BG$ onto a group, between Thomason-fibrant categories, that is not an isomorphism in $\mathrm{Cat}\,\mathrm{hot}$.}
\end{enumerate}

\textit{Proof.}
Let $G_j$ be the group of permutations of $\{0,\dots,j\}$, with $G_j\subset G_{j+1}$ by fixing $j+1$, and $G=\bigcup G_j$. Let $T_n$ have the objects $0,\dots,n$, with $\operatorname{Hom}(i,j)=G_j$ if $i\leq j$ and empty otherwise, and composition by multiplication in the group of the target. The inclusions $T_n\subset T_{n+1}$ are full. The functor $r_n:T_n\to\mathsf BG_n$ that forgets the objects and the functor $j_n:\mathsf BG_n\to T_n$ onto the object $n$ satisfy $r_nj_n=\mathrm{id}$, and the arrows $i\to n$ labelled by the unit form a natural transformation $\mathrm{id}\Rightarrow j_nr_n$. So $T_n\cong\mathsf BG_n$ in $\mathrm{Cat}\,\mathrm{hot}$.

Fibrancy. A category is Thomason fibrant when it has the right lifting property for the functors $c\,\mathrm{Sd}^2\Lambda\to c\,\mathrm{Sd}^2\Delta$ associated with the horn inclusions, and these are sieves \citep[Proposition~4.2]{Thomason1980}. Let $F:A\to T_n$ be given on such a sieve $A\subset B$. The functor $r_nF$ extends to $B$, because the nerve of a groupoid is a Kan complex. Extend $F$ by sending every object of $B$ outside $A$ to $n$ and every arrow that meets the complement of $A$ to the arrow with the label given by the extension of $r_nF$. No arrow leads from the complement into $A$, and $\operatorname{Hom}(i,n)=G_n$ for all $i$, so this is defined, and it is a functor because composition in $T_n$ multiplies labels. So $T_n$ is fibrant, and the union $T$ is fibrant by the preceding commentary on ms.~46.

The shift. For $g\in G_j$ let $\sigma(g)\in G_{j+1}$ fix $0$ and act on $\{1,\dots,j+1\}$ as $g$ does on $\{0,\dots,j\}$; let $s:T\to T$ send $i$ to $i+1$ and the arrow labelled $g$ to the arrow labelled $\sigma(g)$. On $T_n$ both $s$ and the inclusion into $T_{n+1}$ map, by the arrows labelled by the unit, to the functors with the same labels and constant value $n+1$ on objects, and these two are isomorphic by the cycle $c=(0\,1\,\cdots\,n+1)$, since $cgc^{-1}=\sigma(g)$ for $g\in G_n$. This is (a).

The functors $r_n$ are compatible and give $r:T\to\mathsf BG$; the nerve and homotopy groups commute with the union, so $Br$ is a weak equivalence and $BT$ is a $K(G,1)$. Under this identification $s$ induces $\sigma:G\to G$ on $\pi_1$, which is not surjective (its image fixes $0$). This is (b). For (c): if $r$ were invertible in $\mathrm{Cat}\,\mathrm{hot}$, there would be a functor $j:\mathsf BG\to T$ with $rj$ joined to the identity by natural transformations, that is, with $rj:G\to G$ an inner automorphism. But $j$ sends the object to some $i$ and $G$ into $\operatorname{Hom}(i,i)=G_i$, so the image of $rj$ lies in $G_i\neq G$.\hfill$\square$

Consequences for these pages. The assertion of ms.~47, that a morphism of such systems inducing an isomorphism of ind-objects induces a weak equivalence of the unions, is false: $s$ induces the identity of the ind-object. The question of ms.~48, whether two maps of unions that are homotopic on every stage inside a later stage are homotopic, has the answer no, with homotopies of length three at every stage; what fails is their compatibility, not the growth of their length as in the fence example above. The union of fibrant stages is fibrant, but the union $T$ of stages each isomorphic in $\mathrm{Cat}\,\mathrm{hot}$ to a groupoid is not isomorphic to one, by (c). And a weak equivalence between Thomason-fibrant categories need not be an isomorphism in $\mathrm{Cat}\,\mathrm{hot}$ (ms.~47, ``sauf erreur''): fibrancy is not the envelope property.
\end{dcomment}
